\section*{Chapitre \ref{ch:pl:dataframes} --- Dataframes}

\begin{solution}{exo:pl:dataframes:1}
Fermée et étiquetée à gauche~: dans la barre $[10\!:\!05, 10\!:\!10)$, étiquetée 10:05. Fermée et étiquetée à droite~: dans la barre
$(10\!:\!00, 10\!:\!05]$, étiquetée également 10:05. Même étiquette, intervalles différents, distants de cinq minutes.
\end{solution}

\begin{solution}{exo:pl:dataframes:2}
En \omterm{def:pl:dataframes:shape}{format large}~: une ligne par minute (09:31, 09:32), colonnes AAA et BBB. Une minute pendant laquelle BBB n'a fait l'objet d'aucune transaction porte une valeur manquante dans la colonne BBB~;
repasser en \omterm{def:pl:dataframes:shape}{format long} et retirer les valeurs manquantes redonne les quatre lignes d'origine. Une cellule manquante signifie \emph{absence d'observation}, et non un prix nul,
et le code qui la comble doit préciser par quelle valeur.
\end{solution}

\begin{solution}{exo:pl:dataframes:3}
Une seule partition, \texttt{date=5}, et quatre des six colonnes~: \texttt{ts} et \texttt{price}, que la requête renvoie, et \texttt{symbol}
et \texttt{date}, dont le filtre a besoin.
\end{solution}

\begin{solution}{exo:pl:dataframes:4}
La barre fermée à droite $k$ est $(kw, (k+1)w]$. Pour des entiers, $kw < t \le (k+1)w$ équivaut à $kw \le t-1 < (k+1)w$, donc $k = \lfloor (t-1)/w \rfloor$.
La troncature vers zéro ne coïncide avec la partie entière inférieure que pour les nombres positifs ou nuls~: pour $t - 1 = -7$ et $w = 5$, la partie entière inférieure vaut $-2$ et la
troncature $-1$, si bien que les instants juste en dessous de zéro sont rangés dans la barre juste au-dessus, qui devient deux fois plus large. Des instants négatifs apparaissent dès que
les instants sont mesurés à partir d'une référence comme l'ouverture de séance.
\end{solution}

\begin{solution}{exo:pl:dataframes:5}
La série comblée compte quatre rendements de plus, tous nuls, pour les barres vides~; le rendement qui suit chaque trou est le même dans les deux séries, mais dans la
série sans barres vides, il couvre dix minutes et il est traité comme un rendement de cinq minutes. La volatilité glissante sur 20 barres de cette série mélange alors
des rendements de cinq et de dix minutes~: elle surestime la volatilité à cinq minutes autour des trous, et sa fenêtre couvre plus de 100 minutes.
\end{solution}

\begin{solution}{exo:pl:dataframes:6}
Les lectures, les filtres et les agrégations peuvent chacun traiter un lot et le transmettre, en ne gardant qu'un état borné. Le pipeline du chapitre trie les
transactions et les rattache aux cotations par une jointure as-of, ce qui exige que les deux côtés soient triés~: ces opérateurs retiennent toute leur entrée, et le moteur d'\omterm{def:pl:dataframes:outofcore}{exécution en flux}
se replie sur le moteur en mémoire lorsqu'un opérateur n'est pas implémenté en flux.
\end{solution}

\begin{solution}{exo:pl:dataframes:7}
Sur le mois des tests, pour le symbole S003 (7\,960 transactions horodatées à la seconde)~: les trois moteurs renvoient des barres identiques selon chacune des
quatre conventions, et le volume total est le même dans les quatre cas --- chaque transaction appartient à une seule barre~; les conventions déplacent du volume
d'une barre à l'autre et renomment les barres~; elles ne créent ni ne perdent de volume.
\end{solution}

\begin{solution}{exo:pl:dataframes:8}
Trois défauts. Les conventions sont les valeurs par défaut de la bibliothèque et ne sont pas énoncées~: pandas ferme et étiquette les barres à gauche pour la plupart des fréquences, mais
à droite pour les fins de mois, de trimestre, d'année et de semaine. Les rendements sont calculés sur une table triée par horodatage et non regroupée par symbole, de sorte que chaque
rendement compare la barre d'un symbole à la ligne précédente, qui porte la barre d'un autre symbole. Les barres vides sont soit retirées, soit absentes, de sorte que les rendements
couvrent des intervalles inégaux, et les rangs les comparent. Calculer les rendements par symbole, énoncer et fixer les conventions de barres, et décider de ce qu'est une
barre vide avant tout classement.
\end{solution}

\begin{solution}{pb:pl:dataframes:1}
\begin{enumerate}
\item Deux~; les horodatages sont tronqués à la seconde, et une seconde sur 300 est une frontière de cinq minutes, si bien qu'environ 375/300 transactions sont
attendues sur les edges.
\item Quatre intervalles~; 600 actions changent de barre (chaque transaction sur un edge quitte une barre et entre dans sa voisine).
\item Cinq minutes~: la convention à gauche désigne un intervalle par son début, celle à droite par sa fin.
\item Barres vides retirées~: 74 barres et 73 rendements~; barres vides comblées~: 78 barres et 77 rendements.
\item L'une ou l'autre, appliquée partout et consignée avec les données~: largeur, côté fermé, étiquette, traitement des barres vides et règle de départage. Avec
des étiquettes à gauche, la clôture de la barre n'est connue qu'à l'étiquette plus la largeur~: rattacher une barre à d'autres données à son étiquette constitue un look-ahead.
\item Chaque table intermédiaire, en entier, à chaque étape.
\item L'élagage des partitions, le \omterm{def:pl:columnar-formats:pushdown}{report des projections} (projection pushdown) et le \omterm{def:pl:columnar-formats:pushdown}{report des prédicats} (predicate pushdown).
\item Le moteur voit tout le calcul avant d'en exécuter la moindre partie~: il peut retirer les colonnes inutilisées, déplacer les filtres vers la lecture et choisir
les algorithmes.
\item Les tris et la jointure as-of, qui exigent toute leur entrée, dans l'ordre.
\item La division entière de DuckDB tronque vers zéro~; l'indice de barre a besoin de la partie entière inférieure, qui diffère pour les nombres négatifs.
\item Barres d'une minute, rendements, rangs en coupe transversale et rang moyen de chaque symbole~; transactions rattachées par jointure as-of et
fourchette effective moyenne de chaque symbole. Les réponses sont arrondies à neuf décimales et comparées pour l'égalité avant que le moindre temps ne soit consigné.
\item pandas 895~Mo, Polars eager 769, lazy 654, streaming 614, DuckDB 317.
\item Polars streaming est le plus rapide (0.82~s), DuckDB utilise le moins de mémoire (317~Mo).
\item Il charge tout le mois en mémoire, transforme les colonnes de chaînes en objets Python, et chaque étape matérialise une nouvelle table.
\item Les ressources consommées sont conservées à travers \texttt{exec} sous Linux, si bien que le pic de \texttt{getrusage} du processus fils peut inclure la mémoire du processus
dont il est issu.
\item \textbf{Résultat nommé.} Deux conventions de barres sur les mêmes transactions diffèrent pour 4 des 74 volumes de cinq minutes (600 actions) et étiquettent le même
intervalle à cinq minutes d'écart, et le comblement des barres vides fait passer 73 rendements à 77~; le même pipeline de recherche sur un mois exige 895~Mo
hors imports dans pandas, 769, 654 et 614~Mo dans Polars eager, lazy et streaming, et 317~Mo dans DuckDB, de 0.8 à 1.3~s, avec
des réponses identiques.
\item Pour un mois, n'importe lequel. Pour cinq ans (soixante fois les données), un moteur à \omterm{def:pl:dataframes:eager}{évaluation paresseuse} sur les partitions, DuckDB ou Polars streaming, ou le
pipeline exécuté par jour ou par mois lorsque le calcul le permet (les rendements d'un jour à l'autre ont besoin de la dernière barre de la veille).
\item Une bibliothèque commune d'idiomes (barres, jointures as-of, fenêtres, rangs, pivots), avec les conventions en arguments et consignées dans les
métadonnées des jeux de données~: largeur, côté fermé, étiquette, barres vides, règle de départage, unités et fuseau horaire.
\item Contre le changement silencieux d'un résultat par la valeur par défaut d'un moteur --- une règle de départage, un regroupement en classes, un arrondi, le traitement des valeurs manquantes ---
et contre un changement de comportement d'une nouvelle version d'un moteur.
\item Des conventions énoncées, une seule implémentation testée sur tous les moteurs, et un coût mesuré.
\end{enumerate}
\end{solution}

\begin{solution}{iq:pl:dataframes:1}
La largeur, le côté fermé, l'étiquette, le fuseau horaire et les horaires de séance, le traitement des barres vides, la résolution des horodatages (transactions sur les
edges), l'ordre de départage des premier et dernier cours, et les transactions retenues (codes de condition, transactions hors bourse). Puis comparer les deux
ensembles barre par barre sur les mêmes transactions.
\iqlookfor{Conventions d'abord, puis comparaison barre par barre.}
\end{solution}

\begin{solution}{iq:pl:dataframes:2}
L'\omterm{def:pl:dataframes:eager}{évaluation immédiate} exécute chaque opération aussitôt et matérialise chaque table intermédiaire~; l'\omterm{def:pl:dataframes:eager}{évaluation paresseuse} enregistre un plan et l'exécute lorsqu'un
résultat est demandé, après optimisation (reports, élagage, suppression des colonnes inutilisées) et éventuellement en flux. Cela compte lorsque les données
sont volumineuses par rapport à la mémoire ou lues dans de nombreux fichiers, et lorsqu'un pipeline calcule beaucoup de choses qu'il ne renvoie pas.
\iqlookfor{Matérialisation et optimisation.}
\end{solution}

\begin{solution}{iq:pl:dataframes:3}
Barres par symbole (une partition par date et symbole), rendements par symbole, puis un rang par \omterm{def:pl:dataframes:window}{fonction de fenêtre} partitionné par minute, avec une règle de départage énoncée, dans un
moteur paresseux ou SQL sur des partitions par date, un jour à la fois, en conservant la dernière barre de chaque symbole de la veille pour le premier rendement. Environ
$3\,000 \times 390 \times 252 \approx 3\times10^8$ lignes~: une taille pour un moteur qui traite en flux, non pour une table en mémoire et à \omterm{def:pl:dataframes:eager}{évaluation immédiate}.
\iqlookfor{Partitionné, en flux, conventions et règles de départage énoncées.}
\end{solution}

\begin{solution}{iq:pl:dataframes:4}
Ne lire que les colonnes et partitions nécessaires (une lecture paresseuse), utiliser des types compacts (symboles catégoriels, prix et temps entiers), traiter par
jour lorsque le calcul le permet, et basculer le pipeline vers un moteur paresseux ou SQL qui traite en flux~; mesurer le pic avant et après.
\iqlookfor{Projection, types, découpage, moteur paresseux, mesure.}
\end{solution}

\begin{solution}{iq:pl:dataframes:5}
Un contrat par idiome (barres, jointures as-of, fenêtres glissantes, rangs, pivots), avec les conventions en arguments obligatoires~; une implémentation par
moteur derrière lui~; des tests d'équivalence sur des données construites pour toucher les bornes et les cas d'égalité, exécutés en intégration continue à chaque version de moteur~; les
conventions consignées dans les métadonnées de chaque jeu de données~; des benchmarks de temps et de mémoire par moteur pour guider le choix.
\iqlookfor{Conventions explicites et équivalence testée entre moteurs.}
\end{solution}
